\documentclass[12pt,a4paper]{article}
\usepackage{amsmath}
\usepackage{amssymb}
\usepackage{epstopdf}
\usepackage{inputenc}
\usepackage{graphicx}
\usepackage{titletoc} 
\usepackage{fancyhdr}   
\usepackage[a4paper,pdftex]{geometry}	
\usepackage[english]{babel}
\usepackage{xcolor} 
\usepackage{enumerate}
\usepackage{fix-cm} 
\usepackage[notlof]{tocbibind}
\usepackage{amsmath}
\usepackage{listings}
\usepackage{float}
\usepackage{enumitem}
\usepackage{xcolor}
\usepackage{listings}
\definecolor{vgreen}{RGB}{104,180,104}
\definecolor{vblue}{RGB}{49,49,255}
\definecolor{vorange}{RGB}{255,143,102}
\renewcommand\lstlistingname{Appendix}
\renewcommand\lstlistlistingname{Appendix}

\makeatletter
\newcommand*\@lbracket{[}
\newcommand*\@rbracket{]}
\newcommand*\@colon{:}
\newcommand*\colorIndex{%
	\edef\@temp{\the\lst@token}%
	\ifx\@temp\@lbracket \color{black}%
	\else\ifx\@temp\@rbracket \color{black}%
	\else\ifx\@temp\@colon \color{black}%
	\else \color{vorange}%
	\fi\fi\fi
}
\makeatother

\usepackage{trace}

\usepackage{subcaption}
\begin{document}
	\begin{titlepage}
		\begin{center}
			\includegraphics[scale=.4]{Figures/Cover}\\
			\vspace{1cm}
			\bf{ \large {Department of Computer Science and Technology} }
		\end{center}
		
		\vspace{4cm}
		\centering
		\textbf{\Huge Natural Language Processing}
		\vspace{.5cm}
		
		{\Large Assignment 2}

		\vspace{4cm}
		
		\textbf{\LARGE Sahand Sabour}
		
		
		
		\vspace{0.5cm}
		
		{\large 2020280401}
		
		
		\vfill
		
	\end{titlepage}

	\section{Gradient Calculation}
	Assuming that we are given a predicted word vector $v_c$ for the center
word c in skip-gram, and word prediction is made with the following Softmax function:
	\begin{equation}
		y_o = p(o|c) = \frac{exp(u^T_ov_c)}{\sum_{w=1}^{W}exp(u^T_wv_c)}
	\end{equation}
	where w denotes the w-th word and $u_w= (w = 1, ..., W)$ are the “output” word
vectors for all words in the vocabulary. In addition, assuming that the cross entropy loss is used, we would have:
	\begin{equation}
		\mathcal{L} = - \sum_{w=1}^{W} t_i log(y_i) = - log(p(o|c)) = - log (exp(u^T_ov_c)) + log (\sum_{w=1}^{W}exp(u^T_wv_c))
	\end{equation}
	Where t is the label and is either 1 or 0 since the input is one-hot encoded. Therefore, there would only be one element of the sum that is non-zero (for the expected word o). Further simplifying the loss function gives:
	\begin{equation}
		\mathcal{L} = - u^T_ov_c + log (\sum_{w=1}^{W}exp(u^T_wv_c))
	\end{equation}
	Accordingly, we would derive the gradient from this loss as follows:
	\begin{equation}
		\begin{aligned}
			\frac{\delta \mathcal{L}}{\delta v_c} = -\frac{\delta u^T_ov_c}{\delta v_c} + \frac{\delta log(\sum_{w=1}^{W}exp(u^T_wv_c))}{\delta v_c} \qquad \qquad
			\\
			=-u^T_o + (\frac{1}{\sum_{w=1}^{W}exp(u^T_wv_c)})(exp(u^T_wv_c))(u_w^T)\\  =  \sum_{w=1}^{W} p(o|c) u_w^T - u^T_o\qquad \qquad\qquad \qquad \qquad
		\end{aligned}
	\end{equation}
	
	\section{Word2vec Implementation}
	For this assignment, Pytorch was chosen as the deep learning framework for implementing this model. This implementation consists of the following steps:
	\subsection{Data Processing}
	\subsubsection{Reducing Corpus Size}
	We are provided with an unprocessed Wikipedia corpus, which is approximately 10 GB large. Due to the limited computing resources, a smaller version of this file, which included 1/20 of the data and was around 500MB large was created. The newly created corpus was used for the rest of this assignment. 
	
	\subsubsection{Creating Word Count Dictionary}
	For each line in the corpus, there are a number of filters that are used to process the line: first, the line is tokenized using simple processing methods from the gensim library; then, the nltk library is utilized to remove any stop-words; lastly, words that are a set of repeated letters (such as 'aaa', 'aaaa'), which were observed to be fairly frequent in the corpus, are removed. Accordingly, words that pass these set of filters are saved in a dictionary and their frequency is recorded.
	
	\subsubsection{Creating Vocabulary Dictionary}
	Initially, based on the word count dictionary, words that are comparatively less frequent are removed. Accordingly, two dictionaries that give an index to each word and map each word to an index (namely word\_to\_id and id\_to\_word) respectively are created. These will be referred to as vocabulary.
	
	\subsubsection{Sub-Sampling}
	According to the source code of the original word2vec implementation, the probability of a word $w_i$ staying in the vocabulary is calculated as 
	\begin{equation}
		P(w_i) = (\sqrt{\frac{count(w_i)}{0.001}} +1 ) \times \frac{0.001}{count(w_i)}
	\end{equation}
	where count is the number of times a word has been seen in the corpus. According to this equation, words that occur frequently have a higher chance of being kept in the vocabulary and less frequent words are more likely to be removed. The same equation is used in the implementation for this assignment to perform sub-sampling on the corpus.
	
	\subsubsection{Negative Sampling}
	In a skip-gram architecture, for each word in the corpus, referred to as a center word, there are a number of words, referred to as context words, that are shown near this word. The task of this architecture is to find context words given center words. However, it is not satisfactory if all the samples presented to the model are positive: meaning they are all actual context words of the given center word. In order to have meaningful training, we need to present negative samples in which the selected words are not context words for the given center word. For this task, a list of words in the vocabulary is initially created, where each word is repeated as many times as its number of occurrences in the word count dictionary. Accordingly, this list is randomly shuffled and each time a constant number of elements is read from this list as to provide negative samples.
	
	\subsection{Training}
	First, we need two embeddings: namely word and context embeddings. These embeddings are similar in shape (vocabulary size $\times$ embedding dimension) and are both initialized with random values. Upon receiving a new batch of inputs, which consists of center words and their corresponding positive and negative samples, the word embedding is used to look up the center word's word vector while the context is used to collect the vector of the other words. Accordingly, the goal is to maximize the similarity between a center word's vector with the vectors of its context words (positive samples) and minimize the similarity between the word and its negative samples. Since the softmax log likelihood loss is computationally expensive, sigmoid log likelihood loss is utilized in this implementation instead.
	\\
	As required three models with embedding dimensions 100, 200, and 300 were respectively trained. The data was provided to the model in batches of 128, with window size of 5 and 5 negative samples. The model was trained for 5 epochs with the learning rate of 0.001.
	
	\subsection{Evaluation}
	\subsubsection{Spearman's Correlation Coefficient}
	The trained word embeddings were evaluated on the wordsim dataset, which provides a pair of words as well as a mean human-annotated score of their similarity. In this section, the word embeddings of different dimensions were used to find the cosine similarities between the words in each pair and the resulting values were compared to the human annotation results via calculating the spearman's correlation coefficient. This implementation utilized the spearman function from the nltk library and the obtained results are provided respectively in the table below:
	\begin{table}[H]
		\begin{center}
			\begin{tabular}{|c|c|}
				\hline
				Embedding size & Correlation Coefficient (\%)  \\ \hline
				100 &         52.14                \\ \hline
				200 &    56.27     \\ \hline
				300 &    59.49  \\ 
				\hline             
			\end{tabular}
		\end{center}
	\end{table}

\vspace{-0.6cm}
	\noindent Hence, it can be observed that increasing the embedding dimension is rather beneficial for creating more accurate word vectors. This is evident as more dimensions allow for more features to be captured and hence, increased details in modeling each word. However, increasing the dimensions considerably increases the training time and computation complexity. Hence, in practice, 100 is used as the embedding dimension for smaller datasets while 300 is used for large corpora.

	
	\subsubsection{Embedding Visualization}
	In addition, visualizing word embeddings could also be utilized to demonstrate the power of these embeddings since the similarities between different words are highlighted when illustrated. The below figure illustrates visualization of different word embeddings with a list of arbitrary words.
	
	\begin{figure}[H]
		\centering
		\begin{subfigure}[H]{0.32\textwidth}
			\centering
			\includegraphics[width=\textwidth]{Figures/visual_100}
			\caption{$dim=100$}
		\end{subfigure}
		\begin{subfigure}[H]{0.32\textwidth}
			\centering
			\includegraphics[width=\textwidth]{Figures/visual_200}
			\caption{$dim=200$}
		\end{subfigure}
		\begin{subfigure}[H]{0.32\textwidth}
			\centering
			\includegraphics[width=\textwidth]{Figures/visual_300}
			\caption{$dim=300$}
		\end{subfigure}
		\caption{Visualization of different embeddings}
	\end{figure}

	\vspace{-0.5cm}
	\noindent Based on the above figures, it can be observed that, although a slight improvement, in embeddings with higher dimensions the connection between different words of the same category is better realized. For instance, with higher dimensions words relating to big corporations such as apple and amazon move closer together. The same case stands for words such as shanghai, beijing, and china which evidently, are in the same region and therefore, highly related.
	
	\subsubsection{Word Analogy}
	The last task in the evaluation is word analogy. Word analogy can be simplified as the task of correctly finding a word that has the same connection with a given word that another pair of words have. For instance, if the pair of words (France and Paris) are given, the connection between these two words is (country and city). Hence, if we were to give a word that has a similar connection to a country name such as 'England', that word would be 'London': Paris is to France as London is to England. Due to the limited training of the word embeddings, due to limited resources as mentioned, this step of the evaluation was carried fairly briefly by manually testing the output of embeddings in response to different sets of arbitrary pairs. It was observed that all embeddings were able to perform this task to a level: for instance, all the embeddings correctly guessed 'paris' when words ['england', 'london', 'france'] were given. More thorough exploration can be expected from future work.
	
	\section{Word2vec Improvement}
	The approach proposed by [1] is implemented in this assignment to improve the word2vec architecture by adding sense embeddings. The methodology for this addition, as well as evaluation with the base model and discussion of the findings will be provided in this section.

	\subsection{Methodology}
	This method builds upon the skip-gram architecture, which was introduced, implemented and evaluated in the previous section. Hence, we are first required to train the word embeddings on the corpus and obtain trained word vectors for each word. Accordingly, we would initialize and update sense embeddings for words in the corpus based on their gloss. In word knowledge corpora such as WordNet, which are a combination of thesauruses and dictionaries, for each word, different meanings (referred to as gloss) as well as a group of its synonyms for this meaning (referred to as synsets) are provided. For the given input data (i.e. Wikipedia), we first read each line and perform the following process:
	\begin{enumerate}
		\item For each word in the line, we utilize the nltk library to do Part of Speech (POS) tagging. However, only words that are Nouns (N), Verbs (V), Adjectives (J), and Adverbs (R) are recorded and used for further processing.
		\item For each of the selected words, we look at the trained word embeddings and obtain their corresponding word vectors. Accordingly, a dictionary with the word as its key and the vector as its value would be created.
		\item For each word in the dictionary, we look up the list of its different senses. Then, for each sense, we retrieve its corresponding key and gloss. 
		\item Since each gloss is a definition of the sense and therefore a sentence, we would repeat the POS tagging step for each gloss and look up the word vectors for each of the selected words. Words whose vector's cosine similarity to the initial word exceeds an arbitrary threshold would be recorded. Then, the average of the selected word vectors would be recorded as the value for the initial word's each sense vector.
		\item Next, different senses for each words are sorted based on their length.
		\item Following the last step, the context vector would be initialized and by the average of word vectors of the words in the sentence. Then, for each word in the context vector, the cosine similarity between this word and the initial word's word vector is calculated and the context word with the highest similarity is selected.
	\end{enumerate}

	
	\subsection{Evaluation}
	For the evaluation of this method and comparison with the base skip-gram architecture, the SCWS dataset was utilized. Each line in this dataset, gives two words, two corresponding sentences (one sentence for each word) that shows their use case, and a set of human annotated scores for their similarities. The similarity score for the two words was calculated as the mean score of the human annotations. Accordingly, for evaluating the base skip-gram model, similar to wordsim dataset, the cosine similarity between the word vectors of the words is calculated. As for the word sense evaluation, the new model is provided with the two words as well as the sample sentences so that the sense vectors of these words are found; afterwards, the cosine similarity between these two sense vectors are calculated instead. The obtained results are provided in the table below:
	\begin{table}[H]
		\begin{center}
			\begin{tabular}{|c|c|c|}
				\hline
				Model (Embedding Size) & Correlation Coefficient (\%) & Accuracy (\%) \\ \hline
				Skip-gram (300) &          62.56     & 82.55      \\ \hline
				Skip-gram with Word Sense (300) &  44.07    &  77.04       \\ 
				\hline             
			\end{tabular}
		\end{center}
	\end{table}

	\vspace{-0.6cm}
	\noindent As demonstrated in the above table, although the new model has a weaker performance, this is believed to be due to the fact that the word embeddings were trained for few epochs and on a much smaller scale of the dataset than what was provided. Since the base model's embeddings are the foundation of creating sense vectors and obtaining context vectors, having sufficiently trained word vectors is essential. This is supported by the fact that the new model outperformed the base model on a subset of the SCWS dataset that was handpicked based on words that the base model was fairly familiar with. The results of this evaluation are provided respectively in the table below:
	\begin{table}[H]
		\begin{center}
			\begin{tabular}{|c|c|c|}
				\hline
				Model (Embedding Size) & Correlation Coefficient (\%) & Accuracy (\%) \\ \hline
				Skip-gram (300) &         69.20    &  84.51          \\ \hline
				Skip-gram with Word Sense (300) &  74.27    & 84.81  \\ 
				\hline             
			\end{tabular}
		\end{center}
	\end{table}
	\vspace{-0.6cm}
	\noindent Hence, it can be expected that the sense embeddings would outperform the base skip-gram architecture, given that the word embeddings are well-trained.
	
	\newpage
	
	\section{References}
	[1] Xinxiong Chen, Zhiyuan Liu, and Maosong Sun. A unified model for word
	sense representation and disambiguation. In Proceedings of the 2014 Conference on Empirical Methods in Natural Language Processing,
	pages 1025–1035, 2014.
	
	\vspace{0.2cm}
	\noindent [2] Eric H Huang, Richard Socher, Christopher D Manning, and Andrew Y Ng.
	Improving word representations via global context and multiple word prototypes. In Proceedings of the 50th Annual Meeting of the Association for
Computational Linguistics: Long Papers-Volume 1, pages 873–882. Association for Computational Linguistics, 2012.
\end{document}